\chapter{Fundamentos Teóricos}

\section{Scoring basado en perfiles ponderados}

El algoritmo implementado busca determinar una combinación óptima de pesos asociada a tres componentes del perfil analizado:

\begin{itemize}
    \item Componente taxonómica ($T$).
    \item Componente socioeconómica ($S$).
    \item Componente funcional ($F$).
\end{itemize}

Cada perfil ponderado se calcula mediante:

\begin{equation}
P_i = W_1T_i + W_2S_i + W_3F_i
\end{equation}

donde:

\begin{equation}
W_1 + W_2 + W_3 = 1
\end{equation}

y además:

\begin{equation}
W_i \ge 0
\end{equation}

Estas restricciones garantizan que el vector de pesos pertenezca al simplex, permitiendo interpretar cada componente como una proporción relativa de la contribución total.

\section{Cálculo del Score}

Sea:

\[
A \in \mathbb{R}^{M\times N}
\]

la matriz de características observadas.

El score asociado a cada muestra se calcula mediante:

\begin{equation}
Score = A \cdot P
\end{equation}

donde:

\begin{itemize}
    \item $A$ representa los datos de entrada.
    \item $P$ representa el perfil ponderado.
\end{itemize}

Cada elemento del vector \textit{Score} corresponde a la puntuación obtenida por una muestra biológica después de combinar las contribuciones de todos los ítems mediante el vector de pesos.

\section{Evaluación mediante AUC}

La calidad de cada candidato es evaluada mediante el área bajo la curva ROC:

\begin{equation}
AUC = ROC\_AUC(y, Score)
\end{equation}

donde:

\begin{itemize}
    \item $y$ corresponde a las etiquetas reales.
    \item $Score$ representa la puntuación calculada.
\end{itemize}

El valor del AUC se encuentra en el intervalo $[0,1]$. En problemas de clasificación binaria, un valor cercano a $0.5$ indica un comportamiento equivalente al azar, mientras que valores próximos a $1$ representan una excelente capacidad de discriminación entre las clases. En este proyecto se considera de interés práctico el intervalo:

\begin{equation}
0.5 \le AUC \le 1.0
\end{equation}

\section{Función Objetivo}

El problema de optimización consiste en encontrar el vector de pesos $W$ que maximiza la capacidad discriminativa del modelo medida mediante el AUC. Formalmente, el problema se expresa como:

\begin{equation}
\max_{W} AUC(y,\;Score(W))
\end{equation}

donde $Score(W)$ depende del vector de pesos utilizado para construir el perfil ponderado.

Cada candidato representa una posible solución del problema y es evaluado de manera independiente, permitiendo aprovechar diferentes estrategias de paralelización sin modificar la lógica del algoritmo.

\section{Validación del Scoring}

Además de maximizar el AUC, el proyecto considera una medida de consistencia basada en un umbral de decisión $\theta$. Esta métrica cuantifica la proporción de muestras correctamente clasificadas:

\begin{equation}
Consistencia=
\frac{\left|\left\{j:Score_j>\theta,\;y_j=1\right\}\right|}{5}
+
\frac{\left|\left\{j:Score_j\le\theta,\;y_j=0\right\}\right|}{5}
\end{equation}

Un modelo se considera satisfactorio cuando cumple:

\begin{equation}
Consistencia \ge 0.8
\end{equation}

Esta validación complementa el análisis realizado mediante la métrica ROC-AUC, permitiendo verificar que el vector de pesos obtenido genera una separación adecuada entre las muestras sanas y las muestras enfermas.

\section{Búsqueda aleatoria}

Los vectores de pesos candidatos son generados mediante una distribución Dirichlet:

\begin{equation}
W \sim Dirichlet(1,1,1)
\end{equation}

Esta distribución garantiza automáticamente que los pesos satisfacen las restricciones del simplex, es decir, que todos los componentes son no negativos y su suma es igual a uno.

Cada candidato generado es evaluado de forma independiente mediante el cálculo del perfil ponderado, el score correspondiente y la métrica ROC-AUC. Esta independencia entre evaluaciones constituye la principal fuente de paralelismo explotada por las implementaciones desarrolladas en Python, OpenMP, MPI y CUDA.